% titel: Betrachtungen in verschiedenen Systemen: rotierendes Buro
% skript: Kap. Rotierende Bezugssysteme, Corioliskraft

Eine Kugel wird in einem um die eigene Achse rotierenden, kreisrunden Raum
abgeschossen. Je nach Bezugssystem beschreibt man dasselbe physikalische
Ereignis unterschiedlich: im Inertialsystem fliegt die Kugel geradeaus und das
Ziel bewegt sich; im mitrotierenden System krümmt die Kugel ihre Bahn
(Corioliskraft). Beide Betrachtungen müssen dieselbe Abweichung liefern.

\begin{werkzeugkasten}[Werkzeuge]
\textbf{Im Inertialsystem} (ruhend von aussen): Die Kugel hat nur die
Anfangsgeschwindigkeit $\vec v_0$ und die Flugzeit zum Ziel
$t=\text{Distanz}/v_0$. Das Ziel (an der Wand) rotiert mit $\vec\Omega$:
in der Zeit $t$ dreht sich der Raum um $\Omega t$, und der Zielpunkt
verschiebt sich tangential um
\begin{equation}
  \Delta d = r\,\Omega t\;.
\end{equation}

\textbf{Im rotierenden System:} Trägheitskräfte für die Kugel. Bei
(konstant angenommener) Geschwindigkeit $v_0$ wirkt die
Coriolisbeschleunigung
\begin{equation}
  a_C = 2\,\Omega\, v_0
\end{equation}
senkrecht zur Bewegung; Querablenkung (gleichmässig beschleunigt über die
Flugzeit $t$):
\begin{equation}
  \Delta d = \tfrac12 a_C t^2 = \Omega\, v_0\, t^2.
\end{equation}

\textbf{Zentrifugalbeschleunigung} im Abstand $r$ vom Zentrum:
\begin{equation}
  a_\mathrm{zentri} = \Omega^2 r\;.
\end{equation}

\textbf{Konstanten:} $\Omega = 2\pi\cdot 1\,\mathrm{s^{-1}}$ (eine Umdrehung
pro Sekunde), $v_0=\qty{500}{m/s}$, $d=\qty{10}{m}$ (Radius $r=\qty{5}{m}$).
\end{werkzeugkasten}

\begin{werkzeugkasten}[Vorgehen in allgemeiner Form]
Flugzeit $t$ aus Distanz und $v_0$; Ablenkung im Inertialsystem über die
Rotation des Zielpunkts, im rotierenden System über die Coriolisbeschleunigung
bestimmen - beide müssen gleich sein. Für die Beschleunigung am Aufenthaltsort
die Zentrifugalkraft nutzen.
\end{werkzeugkasten}

\begin{achtung}[Typische Fallen]
\begin{itemize}
  \item Die Flugzeit $t$ aus der \emph{im Raum} zurückgelegten Strecke und
       $v_0$; die Rotation ändert diese nur zweiter Ordnung.
  \item In $\Delta d=\tfrac12 a_C t^2$ die Einheit der Zeit (s) korrekt
       einsetzen; $\Omega$ in rad/s.
\end{itemize}
\end{achtung}